% Generated from CANONICAL_CORE.tex lines 988-1018
\section{Counterexamples for Necessity}

\subsection{Without H1 (Convexity)}

Let ${C_{\mathrm{proj}}} = \{x \in \mathbb{R}^2 : \norm{x} = 1\}$ (unit circle, not convex).

For $\Psi = (0, 0)$, every point on ${C_{\mathrm{proj}}}$ is equidistant. Projection is not unique.

\subsection{\texorpdfstring{Without H2 ($\norm{K} < 1$)}{Without H2 (norm K < 1)}}

Take $\Hilb={C_{\mathrm{proj}}}=\mathbb R^2$, $\U=\{0\}$, $\Gamma(0)=0$ and let $K$ be rotation by $\pi/2$. The point $0$ is fixed, but every nonzero orbit cycles with period four and fails to converge. Thus removing strict contraction can destroy convergence even when a unique fixed point remains.

To obtain failure of existence as well, take $K=\operatorname{Id}$ and $\Gamma(0)=v\neq0$ on the same space. The equation $x=x+v$ has no solution. Both examples preserve H1, H3 and H4.

\subsection{Without H3 (Completeness)}

H1 already specifies a Hilbert space, so H3 cannot be removed independently while keeping H1 unchanged. To test the extension to a real inner-product space without completeness, take $c_{00}$, the real finite-support sequences with the $\ell^2$ norm, and ${C_{\mathrm{proj}}}=c_{00}$. Let $S(x_1,x_2,\ldots)=(0,x_1,x_2,\ldots)$, $K=S/2$, $\U=\{0\}$ and $\Gamma(0)=e_1$. Projection onto ${C_{\mathrm{proj}}}$ is the identity and $T(x)=Sx/2+e_1$ is a strict contraction. Its only possible fixed-point coordinates are $(1,1/2,1/4,\ldots)$, which do not belong to $c_{00}$. The iteration from zero is Cauchy but has no limit in the stated space. This isolates the completeness requirement for the generalized contraction argument.

\subsection{Without H4 (Compactness)}

Take $\Hilb={C_{\mathrm{proj}}}=\mathbb R$, $K=0$, $\U=\mathbb R$ and $\Gamma(u)=u$. The input map is uniformly continuous, but the parameter space is not compact. Each transition has the unique fixed point $f_u^\ast=u$, and $\A^\ast=\mathbb R$ is not compact. Thus compactness of the parameter space cannot simply be dropped while retaining the stated compact-image conclusion. It is sufficient, not necessary, since some maps on noncompact parameter spaces have compact image.

\subsection{\texorpdfstring{Why $\Phi$-Regularity Is a Hypothesis}{Why Phi-Regularity Is a Hypothesis}}\label{app:phi-regularity-counterexample}

Let $\mathcal{I}=\mathbb{R}$ with the Euclidean metric, let $\mathcal{E}=(-\pi/2,\pi/2)$ with the Euclidean metric, and define
\[
\Phi(x)=\arctan(x).
\]
The map $\Phi$ is proper as a map into the open interval $\mathcal{E}$, locally Lipschitz, and non-constant on every nonempty ball. However, no constant $C>0$ can satisfy
\[
|\arctan(x)-\arctan(y)|\geq C|x-y|
\qquad\forall x,y\in\mathbb{R}.
\]
Indeed, with $y=x+1$,
\[
|\arctan(x+1)-\arctan(x)|\to0
\quad\text{as }x\to\infty,
\]
while $|x-y|=1$. Thus properness, local Lipschitz regularity, and non-collapse do not imply a global lower Lipschitz bound. Any theorem using such a bound must carry it as an explicit hypothesis or prove stronger coercive hypotheses.

The same obstruction persists with a complete codomain: let $\mathcal I=\mathcal E=\mathbb R$ and $\Phi(x)=x^3$. This is a proper, locally Lipschitz, injective map between complete spaces and is nonconstant on every nonempty ball. If $|x^3-y^3|\geq C|x-y|$ held globally for $C>0$, choosing $y=0$ and $x=h\neq0$ would imply $h^2\geq C$ for all $h$, a contradiction as $h\to0$. Completeness of the codomain does not supply a positive lower Lipschitz constant.

% =============================================================================
% REFERENCES
% =============================================================================
\nocite{banach,kreyszig,rudin,brezis}
